\section{DCN Core Kinematics}
\label{appendix:dcn_coredisperionplots}
We present here the core \vlsr and linewidth of the cores with DCN (3-2) detections for the remaining 13 protoclusters. G338.93 and G333.60 are presented in the main body of the text.
\begin{figure*}
    \centering
    \includegraphics[width=0.97\textwidth]{Appendix/appendix_figures/core_dispersion/G008.67_B6_spw1_dcn_statcont.fits_nonpbcor.fits_moments_core_dispersion_may23_ref_scale.png}
    \includegraphics[width=0.97\textwidth]{Appendix/appendix_figures/core_dispersion/G010.62_B6_spw1_dcn_statcont.fits_nonpbcor.fits_moments_core_dispersion_may23_ref_scale.png}
    \caption{Core \vlsr (left) and DCN (3-2) linewidths (right) estimated from the DCN (3-2) fits to the continuum cores towards the intermediate, and evolved protocluster G008.67 (top) and G010.62 (bottom), respectively. The circles and stars represent the detected cores with mass estimates below and above 8~\msun, respectively, with the colorscale displaying the fitted parameters from the DCN (3-2) fits (left: core \vlsr, right: linewidth). The core \vlsr is the centroid velocity of the DCN (3-2) fit minus the cloud \vlsr (taken as 37~\kms, and -2~\kms for G008.67, and G010.62, respectively). The grey contours are the 4 $\sigma$ level of the DCN (3-2) moment 0 map. In the right panels the positions of cores without without a DCN (3-2) detection are marked with a black cross, and black star for cores with a mass estimate below and above 8\msun, respectively, and green triangles, and stars represent cores with a Complex-type DCN (3-2) spectra, with a mass estimate below and above 8\msun, respectively.}
\end{figure*}

\begin{figure*}
\ContinuedFloat
    \centering
    \includegraphics[width=0.97\textwidth]{Appendix/appendix_figures/core_dispersion/G012.80_B6_spw1_dcn_statcont.fits_nonpbcor.fits_moments_core_dispersion_may23_ref_scale.png}
    \includegraphics[width=0.97\textwidth]{Appendix/appendix_figures/core_dispersion/G327.29_B6_spw1_dcn_statcont.fits_nonpbcor.fits_moments_core_dispersion_may23_ref_scale.png}
    \caption{Continued: core \vlsr (left) and DCN (3-2) linewidths (right) estimated from the DCN (3-2) fits to the continuum cores towards the evolved and young protocluster G012.80 (top) and G327.29 (bottom), respectively. The core \vlsr is the centroid velocity of the DCN (3-2) fit minus the cloud \vlsr (taken as 37~\kms, and -45~\kms for G012.80, and G327.29, respectively).}
\end{figure*}

\begin{figure*}
\ContinuedFloat
    \centering
    \includegraphics[width=1\textwidth]{Appendix/appendix_figures/core_dispersion/G328.25_B6_spw1_dcn_statcont.fits_nonpbcor.fits_moments_core_dispersion_may23_ref_scale.png}
    \includegraphics[width=1\textwidth]{Appendix/appendix_figures/core_dispersion/G337.92_B6_spw1_dcn_statcont.fits_nonpbcor.fits_moments_core_dispersion_may23_ref_scale.png}
    \caption{Continued: core \vlsr (left) and DCN (3-2) linewidths (right) estimated from the DCN (3-2) fits to the continuum cores towards the young protoclusters G328.25 (top) and G337.92 (bottom), respectively. The core \vlsr is the centroid velocity of the DCN (3-2) fit minus the cloud \vlsr (taken as -43~\kms, and -40~\kms for G328.25, and G337.92, respectively).}
\end{figure*}

\begin{figure*}
\ContinuedFloat
    \centering
    \includegraphics[width=1\textwidth]{Appendix/appendix_figures/core_dispersion/G351.77_B6_spw1_dcn_statcont.fits_nonpbcor.fits_moments_core_dispersion_may23_ref_scale.png}
    \includegraphics[width=1\textwidth]{Appendix/appendix_figures/core_dispersion/G353.41_B6_spw1_dcn_statcont.fits_nonpbcor.fits_moments_core_dispersion_may23_ref_scale.png}
    \caption{Core \vlsr (left) and DCN (3-2) linewidths (right) estimated from the DCN (3-2) fits to the continuum cores towards the intermediate protoclusters G351.77 (top) and G353.41 (bottom), respectively. The core \vlsr is the centroid velocity of the DCN (3-2) fit minus the cloud \vlsr (taken as -3~\kms, and -17~\kms for G351.77, and G353.41, respectively).}
\end{figure*}

\begin{figure*}
\ContinuedFloat
    \centering
    \includegraphics[width=1\textwidth]{Appendix/appendix_figures/core_dispersion/W43-MM1_B6_spw1_dcn_statcont.fits_nonpbcor.fits_moments_core_dispersion_may23_ref_scale.png}
    \includegraphics[width=1\textwidth]{Appendix/appendix_figures/core_dispersion/W43-MM2_B6_spw1_dcn_statcont.fits_nonpbcor.fits_moments_core_dispersion_may23_ref_scale.png}
    \caption{Continued: core \vlsr (left) and DCN (3-2) linewidths (right) estimated from the DCN (3-2) fits to the continuum cores towards the young protoclusters W43-MM1 (top) and W43-MM2 (bottom), respectively. The core \vlsr is the centroid velocity of the DCN (3-2) fit minus the cloud \vlsr (taken as 97~\kms, and 91~\kms for W43-MM1, and W43-MM2, respectively).}
\end{figure*}

\begin{figure*}
\ContinuedFloat
\centering
\includegraphics[width=1\textwidth]{Appendix/appendix_figures/core_dispersion/W51-IRS2_B6_spw1_dcn_statcont.fits_nonpbcor.fits_moments_core_dispersion_may23_ref_scale.png}
\caption{Continued: core \vlsr (left) and DCN (3-2) linewidths (right) estimated from the DCN (3-2) fits to the continuum cores towards the evolved protocluster W51-IRS2. The core \vlsr is the centroid velocity of the DCN (3-2) fit minus the cloud \vlsr (taken as 61~\kms for W51-IRS2).}
\end{figure*}

\begin{figure*}
\ContinuedFloat
\centering
\includegraphics[width=1\textwidth]{Appendix/appendix_figures/core_dispersion/W43-MM3_B6_spw1_dcn_statcont.fits_nonpbcor.fits_moments_core_dispersion_may23_ref_scale.png}
\includegraphics[width=1\textwidth]{Appendix/appendix_figures/core_dispersion/W51-E_B6_spw1_dcn_statcont.fits_nonpbcor.fits_moments_core_dispersion_may23_ref_scale.png}
\caption{Continued: core \vlsr (left) and DCN (3-2) linewidths (right) estimated from the DCN (3-2) fits to the continuum cores towards the intermediate protoclusters W43-MM3 (top) and W51-E (bottom), respectively. The core \vlsr is the centroid velocity of the DCN (3-2) fit minus the cloud \vlsr (taken as 55~\kms, and 93~\kms for W43-MM3, and W51-E, respectively).}
\end{figure*}